\subsection{Окружность: точки или таймер}
\label{sec:missions-circle}

Окружность радиуса $r$ в плоскости XY задаётся параметрически:
\[
  x(\theta) = r\cos\theta,\qquad y(\theta) = r\sin\theta,\qquad z = \mathrm{const},
\]
где $\theta$ --- угол в радианах. Если удобнее думать в градусах, переводите: $\theta = \mathrm{deg}\cdot\pi/180$ или \verb|math.rad(deg)|.

Есть два способа пролететь окружность.

\subsubsection*{Способ 1. Много точек}

Разбить окружность на 12 точек через $30^\circ$ и пройти их счётчиком шагов, как любой маршрут. Дрон будет притормаживать в каждой вершине --- получится двенадцатиугольник. Так устроен пример с индикацией в разделе~\ref{sec:missions-architecture}.

\subsubsection*{Способ 2. Таймер двигает цель}

Каждые 0.1~с сдвигать угол на несколько градусов и отправлять новую точку. Дрон не успевает остановиться --- цель всё время впереди --- и летит плавно:

\begin{codefile}{lua/examples/06_missions/missions_circle_timer.lua}
\end{codefile}

\paragraph{Что происходит.}
\begin{itemize}
    \item Таймер стартует только после \verb|TAKEOFF_COMPLETE|: отправлять точки до взлёта нельзя.
    \item Пока дрон кружит, события \verb|POINT_REACHED| приходят постоянно --- и игнорируются: переменная \verb|state| равна \verb|"CIRCLE"|. Здесь впервые одно и то же событие означает разное в зависимости от этапа, поэтому программа хранит \hyperlink{term:state}{состояние} и проверяет его вместе с событием: \verb|event == Ev.POINT_REACHED and state == "RETURN"|. Подробно этот приём разобран в разделе~\ref{sec:missions-architecture}.
    \item После $360^\circ$ таймер останавливает сам себя, переводит автомат в \verb|"RETURN"| и отправляет центр. Следующий \verb|POINT_REACHED| в состоянии \verb|"RETURN"| означает «центр достигнут», и только тогда посадка.
    \item Функция \verb|near| сверяет с этим событием реальное положение дрона (\verb|Sensors.lpsPosition()|, допуск 0.2~м). Событие о последней точке круга могло прийти чуть позже перехода в \verb|"RETURN"|; без проверки программа приняла бы его за «центр» и посадила дрон на окружности. Правило: если событие могло опоздать, проверяй не только событие, но и то, что оно описывает.
    \item Первая строка таймера, \verb|if state ~= "CIRCLE" then return end|, --- защита: тик, уже запланированный в момент \verb|stop()|, может сработать ещё раз (глава~\ref{sec:leds}), и без проверки он отправил бы лишнюю точку.
\end{itemize}

\subsubsection*{Курс по касательной}

Чтобы дрон летел «носом вперёд», курс должен быть касательным к окружности: $\psi = \theta + 90^\circ$. В том же тике, что и точка, обновляется курс:

\begin{codefile}{lua/examples/06_missions/missions_circle_yaw_timer.lua}
\end{codefile}

Здесь круг ограничен временем: $360^\circ / 3^\circ \cdot 0.1$~с $= 12$~с, после чего таймер останавливается, дрон возвращается в центр и садится по событию. Длительность круга --- как раз «пауза по смыслу»: сам автопилот о конце круга не сообщает. Та же проверка \verb|state ~= "CIRCLE"| в начале таймера не даёт запоздалому тику затереть точку возврата в центр.

\begin{taskbox}[Восьмёрка]
Две окружности с центрами $(\pm r, 0)$, проходимые в противоположных направлениях. Подсказка: два состояния \verb|"LEFT"| и \verb|"RIGHT"|, шаг угла разного знака, переход --- когда угол прошёл $360^\circ$.
\end{taskbox}
